% chapters/11/rangkuman.tex
\bagian{Rangkuman}

\begin{rangkumanbox}
\begin{itemize}
  \item Mesin keadaan menghasilkan niat, sedangkan pencarian jalur menghasilkan
        rute. Menghubungkan keduanya hanya menuntut satu tambahan pada keadaan
        mengejar, sedangkan mesin keadaannya tidak berubah.
  \item Hasil pencarian disimpan pada NPC dan dipakai berkali-kali. Pencarian
        diulang menurut selang waktu tertentu, atau ketika jalurnya sudah habis
        dilalui.
  \item Pengikut jalur bergerak menuju titik tengah petak berikutnya, memakai
        perulangan \code{while} dan toleransi jarak. Tanpa keduanya, NPC
        meninggalkan petak yang dilampauinya atau berhenti selamanya.
  \item Selama selang hitung ulang, jalurnya menjadi usang karena pemain
        berpindah. Itulah harga yang dibayar atas penghematan perhitungan.
  \item Beban pencarian jalur berlipat menurut banyaknya NPC, seringnya
        pencarian diulang, dan mahalnya satu kali pencarian.
  \item Menghitung ulang setiap bingkai sangat boros: pada pengukuran bab ini
        tiga kali lipat pekerjaan hanya menghasilkan perbaikan satu piksel.
        Sepertiga pekerjaan sudah cukup.
  \item Selang yang terlalu panjang mulai merugikan, sebab NPC tertinggal
        mengikuti jalur yang menuju tempat pemain berada sebelumnya. Titik
        keseimbangannya berada pada sekitar 0,30 detik.
  \item Apabila seluruh NPC menghitung ulang pada bingkai yang sama, laju
        bingkai menjadi naik turun. Sebarkan waktu perhitungannya dengan
        memberikan awal yang berbeda kepada setiap NPC.
  \item Mutu perilaku dinilai dari tiga segi: kelancaran, daya tanggap, dan
        beban komputasi. Ketiganya saling bertentangan, sehingga titik
        keseimbangannya harus dicari dengan mengukur.
  \item Skenario uji yang berguna dapat diulang, memicu perilaku yang hendak
        dinilai, memuat ukuran yang dapat dibandingkan, dan mewakili keadaan
        pemakaian yang sesungguhnya.
  \item Perbaikan yang baik berdasarkan pengukuran, terarah pada penyebabnya,
        dan disertai cara mengukur keberhasilannya.
  \item Bab ini menutup \CPMK{2} sekaligus menutup seluruh blok Python. Seluruh
        perkakas yang telah dibangun tetap berlaku pada Blok~III dan Blok~IV;
        yang berubah hanyalah letak tanggung jawabnya, dari kode Anda kepada
        mesin permainan.
\end{itemize}
\end{rangkumanbox}
